\chapter{Poster Consistency Ledger}

This ledger is deliberately concise. It records statements that must remain
consistent across the poster, manuscript, code, and oral answers.

\begin{longtable}{p{0.30\linewidth}p{0.61\linewidth}}
\toprule
Topic & Approved statement\\ \midrule
Receiver model & The current implementation is a Cell-Free multi-user MISO downlink with single-antenna UEs and distributed AP transmit arrays.\\
Participation $b_{mp}$ & It authorizes dedicated sensing transmission for target $p$; it does not gate globally cooperative communication.\\
P2 & Covariance lifting introduces rank-one constraints, and SDR relaxes them. P2 does not preserve an explicit rank-one convex constraint.\\
P3 & P3 denotes the convex DC-SCA subproblem solved at each iteration, not the original MINLP.\\
Step 4 & Certified physical recovery is an algorithmic step: project support, fix $b$, re-optimize, extract if needed, and validate. It is not a same-level theoretical P4.\\
SDR & SDR is a lower bound / relaxation reference, never a deployable physical baseline.\\
Fig. 5 & A continuous-phase stabilization diagnostic before discrete recovery.\\
Fig. 8 & A QoS-requirement sensitivity study, not a global Pareto frontier.\\
M12 runtime & A computational cost measurement, not evidence of real-time feasibility.\\
Global optimality & Not claimed. The reported method is a validated heuristic recovery pipeline for a nonconvex MINLP.\\
\bottomrule
\end{longtable}

\section{Poster-specific wording repairs}
Use ``Step 1--4'' consistently in the solution framework. Keep P1, P2, and P3
in the first three card titles to match the derivation. Call the final card
``Certified Physical Recovery,'' not P4. Use ``QoS-requirement sensitivity''
for the Fig.~8 description. Refer to the method as ``Cell-Free ISAC'' and,
when challenged about receiver dimensionality, state ``MU-MISO downlink
equivalent model.''
